% !TeX root = part03.tex

\documentclass[../final.tex]{subfiles}

\begin{document}

% ============================================================
% Практическое занятие 3
% ============================================================

\practice{3}{17.09.2026}{Взаимодействие с веществом: продолжение}


% ============================================================
% Многократное рассеяние
% ============================================================

\section*{Многократное рассеяние}

При прохождении заряженной частицы через вещество
она многократно рассеивается в кулоновских полях атомов вещества.

Каждое отдельное рассеяние обычно происходит на небольшой угол.
Однако большое число последовательных отклонений приводит
к заметному изменению направления движения частицы.

Поскольку отклонения вправо и влево равновероятны,

\begin{equation*}
    \langle\theta\rangle=0.
\end{equation*}

Поэтому характерный угол многократного рассеяния определяется
не средним углом, а среднеквадратичным отклонением

\begin{equation*}
    \sqrt{\langle\theta^2\rangle}.
\end{equation*}


% ------------------------------------------------------------
% Рисунок: многократное рассеяние
% ------------------------------------------------------------

\begin{rbox}[left=18pt,right=18pt,top=14pt,bottom=10pt]{[]}
    \centering

    \begin{tikzpicture}[
        >=Latex,
        every node/.style={text=rtext}
    ]

        % Входящая траектория
        \draw[
            ->,
            draw=rtext,
            line width=1.0pt
        ]
        (-4.2,0) -- (-2.3,0);

        % Слой вещества
        \draw[
            draw=rtextmuted,
            line width=0.8pt
        ]
        (-2.2,-1.7) rectangle (2.0,1.7);

        % Траектория внутри вещества
        \draw[
            draw=rc3,
            line width=1.45pt
        ]
        (-2.2,0)
        -- (-1.55,0.12)
        -- (-0.95,-0.10)
        -- (-0.25,0.20)
        -- (0.35,0.08)
        -- (1.0,0.37)
        -- (2.0,0.55);

        % Выходящая траектория
        \draw[
            ->,
            draw=rc3,
            line width=1.45pt
        ]
        (2.0,0.55) -- (4.15,1.20);

        % Невозмущённое направление
        \draw[
            dashed,
            draw=rtextmuted,
            line width=0.75pt
        ]
        (2.0,0.55) -- (4.15,0.55);

        % Угол
        \draw[
            draw=rtext,
            line width=0.8pt
        ]
        (2.75,0.55)
        arc[
            start angle=0,
            end angle=17,
            radius=0.75
        ];

        \node at (3.05,0.78) {$\theta$};

        \node[
            above,
            font=\small
        ]
        at (-3.25,0)
        {частица};

        \node[
            below,
            font=\small
        ]
        at (-0.10,-1.7)
        {вещество};

    \end{tikzpicture}

    \captionsetup{
        font=footnotesize,
        labelfont=bf,
        skip=5pt
    }

    \captionof{figure}{
        Многократное рассеяние заряженной частицы
        при прохождении через слой вещества.
        Отдельные малые отклонения складываются
        в конечный угол $\theta$.
    }
    \label{fig:multiple_scattering_scheme}
\end{rbox}


Для частицы с зарядом $ze$, проходящей слой вещества,
средний квадрат угла рассеяния можно записать в виде

\begin{equation*}
    \langle\theta^2\rangle
    =
    16\pi
    \frac{N_A}{A}
    z^2
    \frac{
        Z(Z+1)r_e^2
    }{
        (\beta pc)^2
    }
    (m_ec^2)^2
    \ln\left(183Z^{-1/3}\right)X.
\end{equation*}

Здесь

\begin{equation*}
    \beta=\frac{v}{c},
\end{equation*}

$p$ --- импульс налетающей частицы,
$z$ --- её заряд в единицах элементарного заряда,
$Z$ и $A$ --- зарядовое и массовое числа вещества,
$r_e$ --- классический радиус электрона,
а $X$ --- пройденная толщина вещества.

Введём \rtxt{3}{радиационную длину} $X_0$:

\begin{equation*}
    \frac{1}{X_0}
    =
    4\alpha r_e^2
    \frac{N_A}{A}
    Z(Z+1)
    \ln\left(183Z^{-1/3}\right).
\end{equation*}

Тогда выражение для среднего квадрата угла принимает вид

\begin{equation*}
    \langle\theta^2\rangle
    =
    \frac{4\pi}{\alpha}
    (m_ec^2)^2
    \frac{z^2}{(\beta pc)^2}
    \frac{X}{X_0}.
\end{equation*}

После подстановки численных констант получаем
\rtxt{3}{формулу Росси}:

\begin{equation*}
    \boxed{
        \sqrt{\langle\theta^2\rangle}
        =
        \frac{
            13.6\ \text{МэВ}
        }{
            \beta pc
        }
        z
        \sqrt{\frac{X}{X_0}}
    }.
\end{equation*}

Здесь угол измеряется в радианах,
а $pc$ необходимо подставлять в МэВ.

Из формулы видно, что характерный угол рассеяния

\begin{equation*}
    \sqrt{\langle\theta^2\rangle}
    \propto
    \sqrt{\frac{X}{X_0}},
\end{equation*}

то есть возрастает с толщиной вещества,
но уменьшается с ростом импульса частицы.


% ============================================================
% Задача 1
% ============================================================

\section{Задача 1. Многократное рассеяние протона}

Рассмотрим протон с кинетической энергией

\begin{equation*}
    E=5\ \text{МэВ},
\end{equation*}

проходящий слой вещества толщиной $1$ см.

Для расчёта используем отношение

\begin{equation*}
    \frac{X}{X_0}
    =
    \frac{1}{36.08}.
\end{equation*}

Требуется оценить характерный диапазон углов
многократного рассеяния.

Масса протона имеет порядок

\begin{equation*}
    m_pc^2\simeq1\ \text{ГэВ}.
\end{equation*}

Так как

\begin{equation*}
    E\ll m_pc^2,
\end{equation*}

протон можно считать нерелятивистским.

Для него

\begin{equation*}
    E
    =
    \frac{m_pv^2}{2}
    =
    \frac{p^2}{2m_p}.
\end{equation*}

Отсюда

\begin{equation*}
    p=\sqrt{2m_pE},
    \qquad
    v=\sqrt{\frac{2E}{m_p}}.
\end{equation*}

Следовательно,

\begin{equation*}
    \beta
    =
    \frac{v}{c}
    =
    \frac{1}{c}
    \sqrt{\frac{2E}{m_p}}.
\end{equation*}

В формуле Росси входит произведение $\beta pc$.
Запишем

\begin{equation*}
    \beta pc
    =
    \frac{v}{c}pc
    =
    pv.
\end{equation*}

Поскольку

\begin{equation*}
    p=m_pv,
\end{equation*}

получаем

\begin{equation*}
    pv
    =
    m_pv^2
    =
    2E.
\end{equation*}

Таким образом,

\begin{equation*}
    \boxed{
        \beta pc=2E
    }.
\end{equation*}

При $E=5$ МэВ

\begin{equation*}
    \beta pc
    =
    10\ \text{МэВ}.
\end{equation*}

Для протона

\begin{equation*}
    z=1.
\end{equation*}

Тогда

\begin{align*}
    \sqrt{\langle\theta^2\rangle}
    &=
    \frac{13.6}{10}
    \sqrt{\frac{1}{36.08}}
    \\[3pt]
    &\simeq
    0.23.
\end{align*}

Следовательно,

\begin{equation*}
    \boxed{
        \sqrt{\langle\theta^2\rangle}
        \simeq
        0.23\ \text{рад}
    }.
\end{equation*}

Это соответствует примерно

\begin{equation*}
    0.23\ \text{рад}
    \simeq
    13^\circ.
\end{equation*}

Для основной части распределения возьмём диапазон примерно
в три среднеквадратичных отклонения:

\begin{equation*}
    \theta_{\min}\simeq0,
\end{equation*}

\begin{equation*}
    \theta_{\max}
    \simeq
    3\sqrt{\langle\theta^2\rangle}.
\end{equation*}

Поэтому

\begin{equation*}
    \theta_{\max}
    \simeq
    3\cdot0.23
    \simeq
    0.69\ \text{рад}.
\end{equation*}

Окончательно

\begin{equation*}
    \boxed{
        0
        \lesssim
        \theta
        \lesssim
        0.69\ \text{рад}
    }.
\end{equation*}


% ============================================================
% Тормозное излучение
% ============================================================

\section*{Тормозное излучение}

При движении заряженной частицы в кулоновском поле ядра
её траектория искривляется.
Следовательно, частица движется с ускорением
и испускает электромагнитное излучение.

Такое излучение называется
\rtxt{3}{тормозным излучением}.


% ------------------------------------------------------------
% Рисунок: тормозное излучение
% ------------------------------------------------------------

\begin{rbox}[left=18pt,right=18pt,top=14pt,bottom=10pt]{[]}
    \centering

    \begin{tikzpicture}[
        >=Latex,
        every node/.style={text=rtext}
    ]

        % Ядро
        \fill[
            rtext
        ]
        (0,-1.0) circle (2.8pt);

        \node[
            below=4pt
        ]
        at (0,-1.0)
        {$Ze$};

        % Электронная траектория
        \draw[
            ->,
            draw=rc3,
            line width=1.5pt
        ]
        (-4.0,0.65)
        .. controls
            (-1.9,0.65)
            and
            (-0.7,0.55)
        ..
        (0.10,0.20)
        .. controls
            (0.85,-0.10)
            and
            (1.85,-0.15)
        ..
        (3.5,0.25);

        \node[
            above
        ]
        at (-3.25,0.65)
        {$e^-$};

        % Фотон
        \draw[
            draw=rc3,
            line width=1.0pt
        ]
        (0.35,0.12)
        -- (0.65,0.38)
        -- (0.95,0.20)
        -- (1.25,0.46)
        -- (1.55,0.28)
        -- (1.85,0.54)
        -- (2.15,0.36)
        -- (2.45,0.62);

        \draw[
            ->,
            draw=rc3,
            line width=1.0pt
        ]
        (2.45,0.62) -- (2.8,0.75);

        \node[
            above
        ]
        at (2.15,0.62)
        {$\gamma$};

        % Прицельный параметр
        \draw[
            dashed,
            draw=rtextmuted,
            line width=0.75pt
        ]
        (0,-1.0) -- (0,0.42);

        \node[
            left
        ]
        at (0,-0.25)
        {$b$};

    \end{tikzpicture}

    \captionsetup{
        font=footnotesize,
        labelfont=bf,
        skip=5pt
    }

    \captionof{figure}{
        Тормозное излучение электрона
        при отклонении в кулоновском поле ядра.
    }
    \label{fig:bremsstrahlung_practice3}
\end{rbox}


Мощность излучения ускоренно движущегося заряда
пропорциональна квадрату ускорения:

\begin{equation*}
    P\propto a^2.
\end{equation*}

Кулоновская сила имеет порядок

\begin{equation*}
    F_{\text{кул}}
    \sim
    \frac{Ze^2}{b^2}.
\end{equation*}

Следовательно,

\begin{equation*}
    a
    =
    \frac{F_{\text{кул}}}{m}
    \sim
    \frac{Ze^2}{mb^2}.
\end{equation*}

Отсюда

\begin{equation*}
    \boxed{
        P\propto\frac{1}{m^2}
    }.
\end{equation*}

Поэтому тормозное излучение особенно существенно
для лёгких заряженных частиц, прежде всего для электронов.

Для тяжёлых частиц из-за большой массы
тормозными потерями во многих случаях можно пренебречь.


% ------------------------------------------------------------
\subsection*{Изменение энергии электрона}
% ------------------------------------------------------------

Радиационные потери электрона удобно записать через $X_0$:

\begin{equation*}
    -
    \frac{dE}{dx}
    =
    \frac{E}{X_0}.
\end{equation*}

Разделим переменные:

\begin{equation*}
    \frac{dE}{E}
    =
    -
    \frac{dx}{X_0}.
\end{equation*}

Интегрируя,

\begin{equation*}
    \ln E
    =
    -
    \frac{x}{X_0}
    +C.
\end{equation*}

При

\begin{equation*}
    x=0,
    \qquad
    E=E_0,
\end{equation*}

получаем

\begin{equation*}
    C=\ln E_0.
\end{equation*}

Следовательно,

\begin{equation*}
    \ln\frac{E}{E_0}
    =
    -
    \frac{x}{X_0},
\end{equation*}

и

\begin{equation*}
    \boxed{
        E(x)
        =
        E_0
        e^{-x/X_0}
    }.
\end{equation*}

При прохождении одной радиационной длины,

\begin{equation*}
    x=X_0,
\end{equation*}

энергия уменьшается до

\begin{equation*}
    E
    =
    \frac{E_0}{e}.
\end{equation*}


% ------------------------------------------------------------
% Критическая энергия
% ------------------------------------------------------------

\subsection*{Критическая энергия}

При сравнительно небольших энергиях электрона
основную роль играют ионизационные потери.

При больших энергиях начинают преобладать
радиационные потери.

\rtxt{3}{Критическая энергия} $E_{\text{кр}}$
определяется условием равенства этих двух механизмов:

\begin{equation*}
    \left(
        -
        \frac{dE}{dx}
    \right)_{\text{ион}}
    =
    \left(
        -
        \frac{dE}{dx}
    \right)_{\text{торм}}.
\end{equation*}


% ------------------------------------------------------------
% Рисунок: критическая энергия
% ------------------------------------------------------------

\begin{rbox}[left=18pt,right=18pt,top=14pt,bottom=10pt]{[]}
    \centering

    \begin{tikzpicture}[
        >=Latex,
        every node/.style={text=rtext}
    ]

        % Оси
        \draw[
            ->,
            draw=rtext,
            line width=0.9pt
        ]
        (0,0) -- (7.0,0)
        node[right] {$E$};

        \draw[
            ->,
            draw=rtext,
            line width=0.9pt
        ]
        (0,0) -- (0,4.5)
        node[above] {$-\dfrac{dE}{dx}$};

        % Ионизационные потери
        \draw[
            draw=rtextmuted,
            line width=1.2pt
        ]
        (0.5,3.6)
        .. controls
            (1.1,1.9)
            and
            (2.1,1.25)
        ..
        (3.0,1.15)
        .. controls
            (4.3,1.12)
            and
            (5.2,1.35)
        ..
        (6.25,1.65);

        % Радиационные потери
        \draw[
            draw=rc3,
            line width=1.45pt
        ]
        (0.7,0.25) -- (6.2,3.65);

        % Критическая энергия
        \draw[
            dashed,
            draw=rtextmuted,
            line width=0.75pt
        ]
        (2.55,0) -- (2.55,1.4);

        \fill[
            rc3
        ]
        (2.55,1.4) circle (2pt);

        \node[
            below=4pt
        ]
        at (2.55,0)
        {$E_{\text{кр}}$};

        \node[
            right,
            font=\small
        ]
        at (5.3,3.1)
        {радиационные};

        \node[
            right,
            font=\small
        ]
        at (4.5,1.25)
        {ионизационные};

    \end{tikzpicture}

    \captionsetup{
        font=footnotesize,
        labelfont=bf,
        skip=5pt
    }

    \captionof{figure}{
        Качественное сравнение ионизационных
        и радиационных потерь.
        В точке $E_{\text{кр}}$ их вклады совпадают.
    }
    \label{fig:critical_energy}
\end{rbox}


Для электрона используется оценка

\begin{equation*}
    \boxed{
        E_{\text{кр}}^{e}
        \simeq
        \frac{800}{Z}\ \text{МэВ}
    }.
\end{equation*}

Для частицы массы $M$

\begin{equation*}
    \boxed{
        E_{\text{кр}}^{M}
        =
        E_{\text{кр}}^{e}
        \left(
            \frac{M}{m_e}
        \right)^2
    }.
\end{equation*}

Следовательно, для тяжёлых частиц критическая энергия
быстро возрастает как квадрат отношения масс.


% ============================================================
% Задача 2
% ============================================================

\section{Задача 2. Потери энергии электрона и мюона в алюминии}

Рассмотрим слой алюминия толщиной

\begin{equation*}
    l=1\ \text{см}.
\end{equation*}

Начальная энергия частицы

\begin{equation*}
    E_0=10\ \text{ГэВ}.
\end{equation*}

Для алюминия

\begin{equation*}
    Z=13,
    \qquad
    A\simeq27,
    \qquad
    \rho\simeq2.7\ \frac{\text{г}}{\text{см}^3}.
\end{equation*}

Массовая толщина слоя равна

\begin{equation*}
    X
    =
    \rho l
    =
    2.7\ \frac{\text{г}}{\text{см}^2}.
\end{equation*}

Радиационная длина алюминия

\begin{equation*}
    X_0
    \simeq
    24.1\ \frac{\text{г}}{\text{см}^2}.
\end{equation*}


% ------------------------------------------------------------
\subsection*{1. Электрон}
% ------------------------------------------------------------

Оценим критическую энергию электрона:

\begin{equation*}
    E_{\text{кр}}^{e}
    \simeq
    \frac{800}{13}\ \text{МэВ}.
\end{equation*}

В грубой оценке семинара

\begin{equation*}
    E_{\text{кр}}^{e}
    \sim
    80\ \text{МэВ}.
\end{equation*}

При этом

\begin{equation*}
    10\ \text{ГэВ}
    \gg
    E_{\text{кр}}^{e}.
\end{equation*}

Следовательно, для электрона основными являются
радиационные потери.

Используем

\begin{equation*}
    E
    =
    E_0
    e^{-X/X_0}.
\end{equation*}

Подставляем:

\begin{align*}
    E
    &=
    10
    \exp\left(
        -\frac{2.7}{24.1}
    \right)
    \text{ГэВ}
    \\[3pt]
    &\simeq
    9\ \text{ГэВ}.
\end{align*}

Следовательно,

\begin{equation*}
    \Delta E
    =
    E_0-E
    \simeq
    1\ \text{ГэВ}.
\end{equation*}

Относительная потеря энергии составляет

\begin{equation*}
    \boxed{
        \frac{\Delta E}{E_0}
        \sim
        10\%
    }.
\end{equation*}


% ------------------------------------------------------------
\subsection*{2. Мюон}
% ------------------------------------------------------------

Для мюона

\begin{equation*}
    m_\mu c^2
    \simeq
    100\ \text{МэВ}.
\end{equation*}

Критическая энергия равна

\begin{equation*}
    E_{\text{кр}}^\mu
    =
    E_{\text{кр}}^e
    \left(
        \frac{m_\mu}{m_e}
    \right)^2.
\end{equation*}

Используя оценки

\begin{equation*}
    E_{\text{кр}}^e
    \simeq
    80\ \text{МэВ},
    \qquad
    m_\mu c^2
    \simeq
    100\ \text{МэВ},
    \qquad
    m_ec^2
    \simeq
    0.5\ \text{МэВ},
\end{equation*}

получаем

\begin{align*}
    E_{\text{кр}}^\mu
    &\simeq
    80
    \left(
        \frac{100}{0.5}
    \right)^2
    \text{МэВ}
    \\[3pt]
    &\simeq
    3.2\cdot10^3\ \text{ГэВ}.
\end{align*}

Следовательно,

\begin{equation*}
    10\ \text{ГэВ}
    \ll
    E_{\text{кр}}^\mu,
\end{equation*}

и тормозным излучением для мюона можно пренебречь.

Основными остаются ионизационные потери.

Используем приближённую формулу

\begin{equation*}
    -
    \frac{dE}{dx}
    =
    0.3
    \frac{
        Zz^2
    }{
        A\beta^2
    }
    \left[
        11.2
        +
        \ln
        \left(
            \frac{
                (\beta\gamma)^2
            }{
                Z
            }
        \right)
        -
        \beta^2
    \right]
    \frac{
        \text{МэВ}\cdot\text{см}^2
    }{
        \text{г}
    }.
\end{equation*}

Для мюона с энергией $10$ ГэВ

\begin{equation*}
    \gamma
    =
    \frac{E}{m_\mu c^2}
    \simeq
    \frac{10\ \text{ГэВ}}{100\ \text{МэВ}}
    \simeq
    100.
\end{equation*}

При этом

\begin{equation*}
    \beta\simeq1.
\end{equation*}

Для алюминия и $z=1$

\begin{align*}
    -
    \frac{dE}{dx}
    &\simeq
    0.3
    \frac{13}{27}
    \left[
        11.2
        +
        \ln
        \left(
            \frac{10000}{13}
        \right)
        -
        1
    \right]
    \\[3pt]
    &\simeq
    2.5\,
    \frac{
        \text{МэВ}\cdot\text{см}^2
    }{
        \text{г}
    }.
\end{align*}

Тогда

\begin{align*}
    \Delta E
    &=
    \left(
        -
        \frac{dE}{dx}
    \right)X
    \\[3pt]
    &\simeq
    2.5\cdot2.7\ \text{МэВ}
    \\[3pt]
    &\simeq
    6.8\ \text{МэВ}.
\end{align*}

% В рукописной записи стоит 7.5 МэВ.
% Арифметически 2.5 * 2.7 = 6.75 МэВ.

Относительная потеря энергии

\begin{equation*}
    \frac{\Delta E}{E}
    \simeq
    \frac{6.8\ \text{МэВ}}{10^4\ \text{МэВ}}
    \simeq
    6.8\cdot10^{-4}.
\end{equation*}

То есть

\begin{equation*}
    \boxed{
        \frac{\Delta E}{E}
        <0.1\%
    }.
\end{equation*}

Таким образом, в одном и том же слое алюминия
электрон теряет заметную долю энергии за счёт
тормозного излучения, тогда как для мюона
потери остаются преимущественно ионизационными.


% ============================================================
% Черенковское излучение
% ============================================================

\section*{Черенковское излучение}

Если заряженная частица движется в прозрачной среде
со скоростью, превышающей фазовую скорость света в этой среде,

\begin{equation*}
    v>\frac{c}{n},
\end{equation*}

возникает \rtxt{3}{черенковское излучение}.

Излучение формирует конус с углом $\theta$
относительно направления движения частицы.

% ------------------------------------------------------------
% Рисунок: геометрия черенковского излучения
% ------------------------------------------------------------

\begin{rbox}[left=18pt,right=18pt,top=14pt,bottom=10pt]{[]}
    \centering

    \begin{tikzpicture}[
        >=Latex,
        every node/.style={text=rtext}
    ]

        % -----------------------------------------------------
        % Основные точки
        % -----------------------------------------------------

        % A — точка, где частица находилась раньше
        % B — текущее положение частицы
        % C — точка касания волнового фронта с оболочкой конуса

        \coordinate (A) at (-1.4,0);
        \coordinate (B) at (2.6,0);
        \coordinate (C) at (-0.4,1.732);


        % -----------------------------------------------------
        % Направление движения частицы
        % -----------------------------------------------------

        \draw[
            ->,
            draw=rtext,
            line width=1.0pt
        ]
        (-3.6,0) -- (3.4,0);

        \node[
            above=4pt
        ]
        at (-3.0,0)
        {$v$};


        % -----------------------------------------------------
        % Текущее положение частицы
        % -----------------------------------------------------

        \fill[
            rc3
        ]
        (B) circle (2.5pt);


        % -----------------------------------------------------
        % Образующие конуса
        % -----------------------------------------------------

        \draw[
            draw=rc3,
            line width=1.35pt
        ]
        (B) -- (-3.0,3.23);

        \draw[
            draw=rc3,
            line width=1.35pt
        ]
        (B) -- (-3.0,-3.23);


        % -----------------------------------------------------
        % Волновые фронты
        %
        % Каждый фронт построен от положения,
        % в котором частица находилась раньше.
        % Они касаются образующих конуса.
        % -----------------------------------------------------

        \draw[
            draw=rtextmuted,
            line width=0.75pt
        ]
        ({-3.0 + 2.8*cos(-60)}, {2.8*sin(-60)})
        arc[
            start angle=-60,
            end angle=60,
            radius=2.8
        ];

        \draw[
            draw=rtextmuted,
            line width=0.75pt
        ]
        ({-1.4 + 2.0*cos(-60)}, {2.0*sin(-60)})
        arc[
            start angle=-60,
            end angle=60,
            radius=2.0
        ];

        \draw[
            draw=rtextmuted,
            line width=0.75pt
        ]
        ({0.0 + 1.3*cos(-60)}, {1.3*sin(-60)})
        arc[
            start angle=-60,
            end angle=60,
            radius=1.3
        ];

        \draw[
            draw=rtextmuted,
            line width=0.75pt
        ]
        ({1.0 + 0.8*cos(-60)}, {0.8*sin(-60)})
        arc[
            start angle=-60,
            end angle=60,
            radius=0.8
        ];


        % -----------------------------------------------------
        % Радиус волнового фронта AC
        % Это направление распространения света.
        % -----------------------------------------------------

        \draw[
            ->,
            draw=rtext,
            line width=1.05pt
        ]
        (A) -- (C);

        \node[
            left=3pt
        ]
        at (-0.90,0.95)
        {$\dfrac{ct}{n}$};


        % -----------------------------------------------------
        % Путь частицы AB
        % -----------------------------------------------------

        \draw[
            draw=rtextmuted,
            line width=0.8pt
        ]
        (A) -- (B);

        \node[
            below=5pt
        ]
        at (0.60,0)
        {$vt$};


        % -----------------------------------------------------
        % Прямой угол:
        % AC перпендикулярно касательной BC
        % -----------------------------------------------------

        \draw[
            draw=rtext,
            line width=0.8pt
        ]
        (-0.515,1.533)
        -- (-0.342,1.433)
        -- (-0.227,1.632);


        % -----------------------------------------------------
        % Угол Черенкова theta
        %
        % ВАЖНО:
        % theta — угол между направлением движения AB
        % и радиусом волнового фронта AC,
        % а НЕ между осью и образующей конуса.
        % -----------------------------------------------------

        \draw[
            draw=rtext,
            line width=0.85pt
        ]
        (-0.75,0)
        arc[
            start angle=0,
            end angle=60,
            radius=0.65
        ];

        \node
        at (-0.83,0.34)
        {$\theta$};


        % -----------------------------------------------------
        % Обозначения точек
        % -----------------------------------------------------

        \fill[rc3] (A) circle (1.8pt);
        \fill[rc3] (C) circle (1.8pt);

        \node[
            below=5pt
        ]
        at (A)
        {$A$};

        \node[
            below=5pt
        ]
        at (B)
        {$B$};

        \node[
            above left=2pt
        ]
        at (C)
        {$C$};

    \end{tikzpicture}

    \captionsetup{
        font=footnotesize,
        labelfont=bf,
        skip=5pt
    }

    \captionof{figure}{
        Геометрия черенковского излучения.
        За время $t$ частица проходит расстояние $AB=vt$,
        а электромагнитная волна распространяется
        на расстояние $AC=ct/n$.
        Радиус $AC$ перпендикулярен образующей конуса $BC$.
    }
    \label{fig:cherenkov_cone}

\end{rbox}
Из геометрии рисунка

\begin{equation*}
    AB=vt,
    \qquad
    AC=\frac{ct}{n}.
\end{equation*}

Поскольку $AC$ является радиусом волнового фронта
в точке касания,

\begin{equation*}
    AC\perp BC.
\end{equation*}

Поэтому треугольник $ABC$ прямоугольный,
а для угла Черенкова $\theta$ получаем

\begin{equation*}
    \cos\theta
    =
    \frac{AC}{AB}.
\end{equation*}

Подставляя $AC$ и $AB$,

\begin{equation*}
    \cos\theta
    =
    \frac{ct/n}{vt}
    =
    \frac{c}{nv}.
\end{equation*}

Так как

\begin{equation*}
    \beta=\frac{v}{c},
\end{equation*}

окончательно

\begin{equation*}
    \boxed{
        \cos\theta
        =
        \frac{1}{\beta n}
    }.
\end{equation*}

Для существования действительного угла необходимо

\begin{equation*}
    \frac{1}{\beta n}\leq1.
\end{equation*}

Следовательно,

\begin{equation*}
    \beta n\geq1.
\end{equation*}

На пороге возникновения излучения

\begin{equation*}
    \beta n=1,
\end{equation*}

поэтому

\begin{equation*}
    \boxed{
        \beta_{\text{пор}}
        =
        \frac{1}{n}
    }.
\end{equation*}


% ============================================================
% Задача 3
% ============================================================

\section{Задача 3. Порог черенковского излучения}

Сравним пороговую кинетическую энергию электрона
для возникновения черенковского излучения
в алмазе и стекле.

Показатели преломления равны

\begin{equation*}
    n_{\text{алмаз}}
    =
    2.4,
    \qquad
    n_{\text{стекло}}
    =
    2.2.
\end{equation*}

Пороговая скорость определяется условием

\begin{equation*}
    \beta_{\text{пор}}
    =
    \frac{1}{n}.
\end{equation*}

Для алмаза

\begin{equation*}
    \beta_{\text{пор}}^{\text{алм}}
    =
    \frac{1}{2.4}
    =
    \frac{5}{12}.
\end{equation*}

Для стекла

\begin{equation*}
    \beta_{\text{пор}}^{\text{ст}}
    =
    \frac{1}{2.2}
    =
    \frac{5}{11}.
\end{equation*}

Обе скорости имеют порядок

\begin{equation*}
    \beta\sim\frac12,
\end{equation*}

поэтому необходимо использовать
релятивистское выражение для энергии.

Лоренц-фактор

\begin{equation*}
    \gamma
    =
    \frac{1}{
        \sqrt{1-\beta^2}
    }.
\end{equation*}

Полная энергия электрона

\begin{equation*}
    E
    =
    \gamma m_ec^2.
\end{equation*}

Нас интересует кинетическая энергия:

\begin{equation*}
    T
    =
    E-m_ec^2
    =
    (\gamma-1)m_ec^2.
\end{equation*}

В оценке из семинара

\begin{equation*}
    \gamma_{\text{алм}}
    \simeq
    \frac{12}{11},
    \qquad
    \gamma_{\text{ст}}
    \simeq
    \frac{11}{10}.
\end{equation*}

Используя

\begin{equation*}
    m_ec^2=511\ \text{кэВ},
\end{equation*}

для алмаза получаем

\begin{align*}
    T_{\text{алм}}
    &\simeq
    \left(
        \frac{12}{11}-1
    \right)
    511\ \text{кэВ}
    \\[3pt]
    &=
    \frac{511}{11}\ \text{кэВ}
    \\[3pt]
    &\simeq
    46\ \text{кэВ}.
\end{align*}

Для стекла

\begin{align*}
    T_{\text{ст}}
    &\simeq
    \left(
        \frac{11}{10}-1
    \right)
    511\ \text{кэВ}
    \\[3pt]
    &\simeq
    51\ \text{кэВ}.
\end{align*}

Следовательно,

\begin{equation*}
    \boxed{
        T_{\text{алм}}
        <
        T_{\text{ст}}
    }.
\end{equation*}

Это связано с тем, что

\begin{equation*}
    n_{\text{алмаз}}
    >
    n_{\text{стекло}},
\end{equation*}

а значит пороговая скорость в алмазе меньше.


% ============================================================
% Взаимодействие гамма-квантов с веществом
% ============================================================

\section*{Взаимодействие $\gamma$-квантов с веществом}

Рассмотрим два механизма взаимодействия
$\gamma$-квантов с веществом:

\begin{itemize}
    \item фотоэффект;
    \item эффект Комптона.
\end{itemize}


% ============================================================
% Фотоэффект
% ============================================================

\subsection*{Фотоэффект}

При фотоэффекте $\gamma$-квант полностью поглощается атомом,
а один из связанных электронов покидает атомную оболочку:

\begin{equation*}
    \gamma+\text{атом}
    \longrightarrow
    \text{атом}^{+}+e^-.
\end{equation*}

Основной вклад в сечение фотоэффекта дают электроны
внутренних оболочек, прежде всего $K$-оболочки.

Полное сечение можно оценить через сечение на $K$-оболочке:

\begin{equation*}
    \boxed{
        \sigma_{\text{ph}}
        \simeq
        \frac{5}{4}
        \left(
            \sigma_{\text{ph}}
        \right)_K
    }.
\end{equation*}

При

\begin{equation*}
    E_K
    \ll
    E_\gamma
    \ll
    m_ec^2
\end{equation*}

имеем

\begin{equation*}
    \boxed{
        \left(
            \sigma_{\text{ph}}
        \right)_K
        \simeq
        10^{-33}
        \frac{
            Z^5
        }{
            E_\gamma^{3.5}
        }
        \ \text{см}^2
    }.
\end{equation*}

Следовательно,

\begin{equation*}
    \sigma_{\text{ph}}
    \propto
    Z^5,
\end{equation*}

а по энергии

\begin{equation*}
    \sigma_{\text{ph}}
    \propto
    E_\gamma^{-3.5}.
\end{equation*}

При больших энергиях,

\begin{equation*}
    E_\gamma
    \gg
    m_ec^2,
\end{equation*}

зависимость от энергии становится более слабой:

\begin{equation*}
    \boxed{
        \left(
            \sigma_{\text{ph}}
        \right)_K
        \simeq
        1.4\cdot10^{-33}
        \frac{
            Z^5
        }{
            E_\gamma
        }
    }.
\end{equation*}

Энергия связи электрона на $K$-оболочке
оценивается как

\begin{equation*}
    \boxed{
        E_K
        \simeq
        14Z^2\ \text{эВ}
    }.
\end{equation*}


% ------------------------------------------------------------
% Качественный график фотоэффекта
% ------------------------------------------------------------

\begin{rbox}[left=18pt,right=18pt,top=14pt,bottom=10pt]{[]}
    \centering

    \begin{tikzpicture}[
        >=Latex,
        every node/.style={text=rtext}
    ]

        \draw[
            ->,
            draw=rtext,
            line width=0.9pt
        ]
        (0,0) -- (6.5,0)
        node[right] {$E_\gamma$};

        \draw[
            ->,
            draw=rtext,
            line width=0.9pt
        ]
        (0,0) -- (0,4.0)
        node[above] {$\sigma_{\text{ph}}$};

        \draw[
            draw=rc3,
            line width=1.5pt
        ]
        (0.45,3.55)
        .. controls
            (0.75,2.3)
            and
            (1.6,1.2)
        ..
        (2.7,0.72)
        .. controls
            (3.8,0.38)
            and
            (4.9,0.23)
        ..
        (6.0,0.15);

    \end{tikzpicture}

    \captionsetup{
        font=footnotesize,
        labelfont=bf,
        skip=5pt
    }

    \captionof{figure}{
        Качественная зависимость сечения фотоэффекта
        от энергии $\gamma$-кванта.
    }
    \label{fig:photoeffect_cross_section}
\end{rbox}


% ============================================================
% Эффект Комптона
% ============================================================

\subsection*{Эффект Комптона}

При комптоновском рассеянии $\gamma$-квант
не поглощается полностью, а передаёт электрону
лишь часть своей энергии:

\begin{equation*}
    \gamma+e^-
    \longrightarrow
    \gamma'+e'.
\end{equation*}

Для одного электрона при малых энергиях,

\begin{equation*}
    E_\gamma
    \ll
    m_ec^2,
\end{equation*}

сечение переходит в томсоновское:

\begin{equation*}
    \boxed{
        \sigma_C^e
        =
        \frac{8}{3}
        \pi r_e^2
    }.
\end{equation*}

При больших энергиях,

\begin{equation*}
    E_\gamma
    \gg
    m_ec^2,
\end{equation*}

получаем

\begin{equation*}
    \boxed{
        \sigma_C^e
        =
        \pi r_e^2
        \frac{
            m_ec^2
        }{
            E_\gamma
        }
        \left[
            \ln
            \left(
                \frac{
                    2E_\gamma
                }{
                    m_ec^2
                }
            \right)
            +
            \frac12
        \right]
    }.
\end{equation*}

То есть при больших энергиях сечение
приближённо уменьшается как

\begin{equation*}
    \sigma_C^e
    \sim
    \frac{
        \ln E_\gamma
    }{
        E_\gamma
    }.
\end{equation*}

В атоме имеется $Z$ электронов,
поэтому сечение на атоме равно

\begin{equation*}
    \boxed{
        \sigma_C
        =
        Z\sigma_C^e
    }.
\end{equation*}


% ------------------------------------------------------------
% Качественный график Комптона
% ------------------------------------------------------------

\begin{rbox}[left=18pt,right=18pt,top=14pt,bottom=10pt]{[]}
    \centering

    \begin{tikzpicture}[
        >=Latex,
        every node/.style={text=rtext}
    ]

        \draw[
            ->,
            draw=rtext,
            line width=0.9pt
        ]
        (0,0) -- (6.5,0)
        node[right] {$E_\gamma$};

        \draw[
            ->,
            draw=rtext,
            line width=0.9pt
        ]
        (0,0) -- (0,3.8)
        node[above] {$\sigma_C$};

        % Низкоэнергетическое плато
        \draw[
            draw=rc3,
            line width=1.5pt
        ]
        (0.45,2.55)
        .. controls
            (1.4,2.55)
            and
            (2.1,2.52)
        ..
        (2.75,2.45)
        .. controls
            (3.35,2.28)
            and
            (3.65,1.70)
        ..
        (4.15,1.12)
        .. controls
            (4.7,0.63)
            and
            (5.3,0.36)
        ..
        (6.0,0.22);

    \end{tikzpicture}

    \captionsetup{
        font=footnotesize,
        labelfont=bf,
        skip=5pt
    }

    \captionof{figure}{
        Качественная энергетическая зависимость
        сечения комптоновского рассеяния.
    }
    \label{fig:compton_cross_section}
\end{rbox}


% ============================================================
% Задача 4
% ============================================================

\section{Задача 4. Кинематика комптоновского рассеяния}

Рассмотрим процесс

\begin{equation*}
    \gamma+e^-
    \longrightarrow
    \gamma'+e'.
\end{equation*}

Начальный электрон покоится.

Пусть исходный фотон имеет энергию $E_\gamma$,
а рассеянный фотон --- энергию $E_\gamma'$.
Угол между направлениями исходного и рассеянного фотонов
обозначим через $\theta$.

Требуется найти

\begin{equation*}
    E_\gamma'
    =
    E_\gamma'(E_\gamma,\theta).
\end{equation*}


% ------------------------------------------------------------
% Рисунок: эффект Комптона
% ------------------------------------------------------------

\begin{rbox}[left=18pt,right=18pt,top=14pt,bottom=10pt]{[]}
    \centering

    \begin{tikzpicture}[
        >=Latex,
        every node/.style={text=rtext}
    ]

        \coordinate (O) at (0,0);

        % Начальный фотон
        \draw[
            draw=rc3,
            line width=1.0pt
        ]
        (-4.0,0)
        -- (-3.65,0.17)
        -- (-3.30,-0.17)
        -- (-2.95,0.17)
        -- (-2.60,-0.17)
        -- (-2.25,0.17)
        -- (-1.90,-0.17)
        -- (-1.55,0.17)
        -- (-1.20,-0.17)
        -- (-0.85,0.17)
        -- (-0.50,0);

        \draw[
            ->,
            draw=rc3,
            line width=1.0pt
        ]
        (-0.50,0) -- (O);

        \node[
            above
        ]
        at (-2.3,0.25)
        {$\gamma,\ E_\gamma$};

        % Начальный электрон
        \fill[
            rtext
        ]
        (O) circle (2.6pt);

        \node[
            below=5pt
        ]
        at (O)
        {$e^-$};

        % Рассеянный фотон
        \draw[
            draw=rc3,
            line width=1.0pt
        ]
        (0,0)
        -- (0.32,0.30)
        -- (0.58,0.20)
        -- (0.90,0.50)
        -- (1.16,0.40)
        -- (1.48,0.70)
        -- (1.74,0.60)
        -- (2.06,0.90)
        -- (2.32,0.80);

        \draw[
            ->,
            draw=rc3,
            line width=1.0pt
        ]
        (2.32,0.80) -- (2.75,1.02);

        \node[
            above
        ]
        at (2.0,0.95)
        {$\gamma',\ E_\gamma'$};

        % Электрон отдачи
        \draw[
            ->,
            draw=rtext,
            line width=1.1pt
        ]
        (O) -- (2.5,-1.35);

        \node[
            below
        ]
        at (2.0,-1.10)
        {$e'$};

        % Исходное направление
        \draw[
            dashed,
            draw=rtextmuted,
            line width=0.7pt
        ]
        (O) -- (3.1,0);

        % Угол theta
        \draw[
            draw=rtext,
            line width=0.8pt
        ]
        (0.85,0)
        arc[
            start angle=0,
            end angle=22,
            radius=0.85
        ];

        \node at (1.15,0.22) {$\theta$};

    \end{tikzpicture}

    \captionsetup{
        font=footnotesize,
        labelfont=bf,
        skip=5pt
    }

    \captionof{figure}{
        Кинематика комптоновского рассеяния.
        Начальный электрон покоится,
        а рассеянный фотон вылетает под углом $\theta$.
    }
    \label{fig:compton_kinematics}
\end{rbox}


Запишем закон сохранения четырёхимпульса:

\begin{equation*}
    P_\gamma+P_e
    =
    P_\gamma'+P_e'.
\end{equation*}

Отсюда

\begin{equation*}
    P_e'
    =
    P_\gamma+P_e-P_\gamma'.
\end{equation*}

В лабораторной системе отсчёта начальный электрон покоится:

\begin{equation*}
    P_e
    =
    \left(
        m_ec,
        \vec 0
    \right).
\end{equation*}

Для фотонов

\begin{equation*}
    P_\gamma^2
    =
    P_\gamma'^2
    =
    0.
\end{equation*}

Для конечного электрона

\begin{equation*}
    P_e'^2
    =
    m_e^2c^2.
\end{equation*}

Поэтому возведём равенство

\begin{equation*}
    P_e'
    =
    P_\gamma+P_e-P_\gamma'
\end{equation*}

в квадрат:

\begin{equation*}
    P_e'^2
    =
    \left(
        P_\gamma+P_e-P_\gamma'
    \right)^2.
\end{equation*}

Получаем

\begin{align*}
    m_e^2c^2
    &=
    P_\gamma^2
    +
    P_e^2
    +
    P_\gamma'^2
    +
    2P_\gamma P_e
    -
    2P_\gamma P_\gamma'
    -
    2P_eP_\gamma'
    \\[3pt]
    &=
    m_e^2c^2
    +
    2P_\gamma P_e
    -
    2P_\gamma P_\gamma'
    -
    2P_eP_\gamma'.
\end{align*}

Сокращая $m_e^2c^2$, имеем

\begin{equation*}
    P_\gamma P_e
    -
    P_\gamma P_\gamma'
    -
    P_eP_\gamma'
    =
    0.
\end{equation*}

Поскольку начальный электрон покоится,

\begin{equation*}
    P_\gamma P_e
    =
    m_eE_\gamma,
\end{equation*}

и

\begin{equation*}
    P_eP_\gamma'
    =
    m_eE_\gamma'.
\end{equation*}

Скалярное произведение четырёхимпульсов
двух фотонов равно

\begin{equation*}
    P_\gamma P_\gamma'
    =
    \frac{
        E_\gamma E_\gamma'
    }{
        c^2
    }
    (1-\cos\theta).
\end{equation*}

Следовательно,

\begin{equation*}
    m_eE_\gamma
    -
    m_eE_\gamma'
    -
    \frac{
        E_\gamma E_\gamma'
    }{
        c^2
    }
    (1-\cos\theta)
    =
    0.
\end{equation*}

Умножая на $c^2$, получаем

\begin{equation*}
    m_ec^2
    \left(
        E_\gamma-E_\gamma'
    \right)
    =
    E_\gamma E_\gamma'
    (1-\cos\theta).
\end{equation*}

Перенесём члены с $E_\gamma'$ в правую часть:

\begin{equation*}
    m_ec^2E_\gamma
    =
    E_\gamma'
    \left[
        m_ec^2
        +
        E_\gamma
        (1-\cos\theta)
    \right].
\end{equation*}

Отсюда

\begin{equation*}
    \boxed{
        E_\gamma'
        =
        \frac{
            E_\gamma
        }{
            1+
            \dfrac{
                E_\gamma
            }{
                m_ec^2
            }
            (1-\cos\theta)
        }
    }.
\end{equation*}

Это формула энергии рассеянного фотона
в эффекте Комптона.

При рассеянии строго вперёд

\begin{equation*}
    \theta=0,
\end{equation*}

и поэтому

\begin{equation*}
    E_\gamma'
    =
    E_\gamma.
\end{equation*}

При обратном рассеянии

\begin{equation*}
    \theta=\pi,
\end{equation*}

поэтому

\begin{equation*}
    1-\cos\theta=2,
\end{equation*}

и минимальная энергия рассеянного фотона равна

\begin{equation*}
    \boxed{
        E_{\gamma,\min}'
        =
        \frac{
            E_\gamma
        }{
            1+
            \dfrac{
                2E_\gamma
            }{
                m_ec^2
            }
        }
    }.
\end{equation*}

По закону сохранения энергии

\begin{equation*}
    E_\gamma+m_ec^2
    =
    E_\gamma'+E_e'.
\end{equation*}

Следовательно,

\begin{equation*}
    E_e'
    =
    E_\gamma
    +
    m_ec^2
    -
    E_\gamma'.
\end{equation*}

Кинетическая энергия электрона отдачи равна

\begin{equation*}
    T_e
    =
    E_e'-m_ec^2.
\end{equation*}

Поэтому

\begin{equation*}
    T_e
    =
    E_\gamma-E_\gamma'.
\end{equation*}

Подставляя найденное выражение для энергии фотона,

\begin{equation*}
    \boxed{
        T_e
        =
        E_\gamma
        \left[
            1-
            \frac{
                1
            }{
                1+
                \dfrac{
                    E_\gamma
                }{
                    m_ec^2
                }
                (1-\cos\theta)
            }
        \right]
    }.
\end{equation*}


\end{document}